% Abhakseite „Das kann ich“ – Zone nur im Gesamt (\mitzone)
%% TODO Ich-kann-Sätze: Platzhalter wie in den Aufgabendateien
\begin{abhakseite}
\ifdefined\mitzone
\abhakgruppe{Kennst du schon}
\abhak{1}{Steigung einer Geraden als Quotient aus Höhenunterschied und waagerechtem Unterschied, Steigungsdreieck lesen und zeichnen, Geradengleichung aus Steigung und Punkt (Punkt-Steigungs-Form)}
\abhak{2}{Funktionswerte berechnen, auch mit e-Term und Bruch, und Werte an Graph und Tabelle ablesen}
\abhak{3}{Einheiten und Zeiten: Uhrzeiten in Stunden seit einem Bezugspunkt und Dezimalstunden in Stunden und Minuten umrechnen, Einheiten von f und x zur Einheit „je“ zusammensetzen (Stück je Tag, Meter je Minute), prozentuale Abweichung als Quotient}
\abhak{4}{Ableitungsfunktion bilden: Potenz-, Faktor- und Summenregel für ganzrationale Funktionen; GK und LK zusätzlich Produkt- und Kettenregel für Produkte aus Polynom und e-Funktion}
\abhak{5}{Gleichungen lösen: linear, quadratisch mit Lösungsformel, Ausklammern und Nullprodukt; GK zusätzlich e-Faktoren als nullstellenfrei erkennen und einfache Exponentialgleichungen logarithmieren; mit dem Rechner numerisch lösen, wenn Teil B es zulässt}
\abhak{6}{Verlauf eines Graphen qualitativ lesen: steigt, fällt, wird steiler oder flacher, Wendestelle als steilste Stelle, Krümmung}
\abhak{7}{Extrempunkte über notwendige und hinreichende Bedingung, Randvergleich auf einem Intervall}
\abhak{8}{Steigung einer Geraden als Quotient aus Höhenunterschied und waagerechtem Unterschied, Steigungsdreieck lesen und zeichnen, Geradengleichung aus Steigung und Punkt (Punkt-Steigungs-Form) – Fehler finden}
\abhak{9}{Steigung einer Geraden als Quotient aus Höhenunterschied und waagerechtem Unterschied, Steigungsdreieck lesen und zeichnen, Geradengleichung aus Steigung und Punkt (Punkt-Steigungs-Form) – selbst rechnen}
\fi
\abhakgruppe{Einheit 1 · Mittlere Änderungsrate und Sekante}
\abhak{10}{Mittlere Änderungsrate}
\abhak{11}{Mittlere Änderungsrate – Fehler finden, Begründen, Anwendung, Darstellungswechsel}
\abhakgruppe{Einheit 2 · Ableitung an einer Stelle}
\abhak{12}{Ableitung an einer Stelle}
\abhak{13}{Momentane Änderungsrate am Graphen über eine eingezeichnete Tangente bestimmen}
\abhak{14}{Graph der Füllhöhe in Abhängigkeit von der Zeit über die Gefäßform auswählen und begründen}
\abhak{15}{Ableitung an einer Stelle – Fehler finden, Begründen, Anwendung, Darstellungswechsel}
\abhakgruppe{Einheit 3 · Von der Sekante zur Tangente}
\abhak{16}{Von der Sekante zur Tangente}
\abhak{17}{Differenzen- und Differentialquotient durch Sekante und Tangente veranschaulichen und im Sachzusammenhang deuten}
\abhak{18}{Von der Sekante zur Tangente – Fehler finden, Begründen, Anwendung, Darstellungswechsel}
\abhakgruppe{Einheit 4 · Die Rate als Funktion}
\abhak{19}{Die Rate als Funktion}
\abhak{20}{Die Rate als Funktion – Fehler finden, Begründen, Anwendung, Darstellungswechsel}
\end{abhakseite}
